%*********************************************************************%
%                                                                     %
%                             Active-DVI                              %
%                                                                     %
%                   Projet Cristal, INRIA Rocquencourt                %
%                                                                     %
%  Copyright 2002,2021 Institut National de Recherche en Informatique %
%  et en Automatique.  All rights reserved. This file is distributed  %
%  under the terms of the GNU Lesser General Public License.          %
%                                                                     %
%  Jun Furuse, Didier R�my and Pierre Weis.                           %
%  Contributions by Roberto Di Cosmo, Didier Le Botlan,               %
%  Xavier Leroy, and Alan Schmitt.                                    %
%                                                                     %
%  Based on Mldvi by Alexandre Miquel.                                %
%*********************************************************************%

\documentclass[12pt]{article}

\usepackage {manual}
\usepackage {fullpage}

\providecommand{\advidriver}{ignore}
\usepackage {graphicx}
\usepackage [\advidriver]{advi}
\usepackage {advi-annot}
\usepackage {bubble}

\usepackage {xspace}
\def \ActiveDVI {Active-DVI\xspace}
\def \advi {\texttt{advi}\xspace}
\def \latex {\LaTeX\xspace}

% For splash inclusion
\usepackage{color}
\definecolor{red}{rgb}{1.0,0.0,0.0}
\definecolor{blue}{rgb}{0.0,0.0,1.0}
\usepackage{graphicx}
\usepackage{tabularx}
%\usepackage{pst-all}
\usepackage{xwindows-colors}
\usepackage[\driver]{hyperref}



% End of splash specific stuff

% We use Hevea to create the HTML version
\usepackage{hevea}

\ifhevea\title{
{\bf\ActiveDVI}\\
\imgsrc{bar.gif}\\
Reference manual\\
{\small Version {\adviversion}}
}
\else
\newdimen \tmpdim
\title {\setbox0 \hbox {\bf\ActiveDVI}\global \tmpdim \wd0
\advance \tmpdim by 4em
\begin{tabular}{c}
\box0 \\
\includegraphics[width=\the\tmpdim]{bar.eps}\\
Reference manual\\
{\small Version {\adviversion}}
\end{tabular}
}
\fi
\author {Didier R{\'{e}}my and Pierre Weis}

\thispagestyle{empty}

\usepackage{makeidx}
\makeindex

\begin{document}

\maketitle

\ifhevea
\renewenvironment{abstract}{\begin{quote}}{\end{quote}}
\else
\def \abstractname {}
\fi

\begin{abstract}
{\bf\ActiveDVI} is a viewer for DVI files that also recognizes a new class of 
\docdef\special's targeted to presentations via laptop computers:
various visual effects can easily be incorporated to the presentation,
via a companion \doctt{advi.sty} {\latex} package. 

\ifhevea\else\bigskip\fi
{\ActiveDVI} is copyright \copyright 2001--{\adviyear},2021 INRIA
and distributed under the Gnu Library General Public License ---see the LGPL
file included in the distribution.

\ifhevea
This manual is  available as
\ahref {manualxxx.html}{several HTML files}
or in \ahref {manual.pdf}{PDF format}.
\fi
\end{abstract}

\vfill
\section*{Acknowledgments and contributors}

{\ActiveDVI} is based on Mldvi, written and 
\ahref{http://www.pps.jussieu.fr/~miquel/soft.html}{distributed}
by \ahref{http://www.pps.jussieu.fr/~miquel/}{Alexandre~Miquel}
under the LGPL license\footnote {At the url
\url{http://www.pps.jussieu.fr/~miquel/soft.html}.}.
This constitutes the core rendering engine.  
{\ActiveDVI} has then been developed by 
%
\ahref{http://jun.furuse.info/}{Jun Furuse},
\ahref{http://gallium.inria.fr/\home{remy}/}{Didier R{\'e}my}, and
\ahref{http://gallium.inria.fr/\home{weis}/}{Pierre Weis},
with contributions by 
\ahref{http://wwwdgeinew.insa-toulouse.fr/\home{lebotlan}/}{Didier Le Botlan},
\ahref{http://www.dicosmo.org/}{Roberto Di Cosmo},
\ahref{http://gallium.inria.fr/\home{xleroy}/}{Xavier Leroy},
\ahref{http://www.pps.jussieu.fr/\home{miquel}/}{Alexandre Miquel}, and
\ahref{http://sardes.inrialpes.fr/\home{aschmitt}/}{Alan Schmitt}.

\clearpage
%\advitoc
\hypertarget{advitoc.first}{}
\tableofcontents
\hypertarget{advitoc.last}{}
%\endadvitoc
\clearpage

\section{Installation}

\ActiveDVI may be available as a package of your linux or macOS distribution.

Otherwise, starting with version 2.0.0 you need an \verb"opam"-based
installation of \verb"ocaml" (see \url {https://opam.ocaml.org/} and \url
{https://ocaml.org/index.fr.html}).

\subsection{Installation of the \advi executable}

The simpliest is to run:
\begin{verbatim}
        opam install advi
\end{verbatim}
this should tell you and install the missing dependencies. 

You may also retreive the github sources from \url
{https://github.com/diremy/advi} either by cloning sources or by retreiving
a tar ball.  In both cases, you must ensure that the following \verb"opam"
packages
\begin{verbatim}
        "dune" {>= "2.5"}
        "graphics" {>= "5.1.1"}
        "camlimages" {>= "5.0.4"}
\end{verbatim}
are installed.  Then just run in the root distribution directory:
\begin{verbatim}
        make all
        make install
\end{verbatim}
This will only install the \verb"advi" executable in the opam hierarchy as
well as a command \verb"advi-latex-files" which you must run as described
below. 

\subsection{Installation of the \advi \latex source files}

As explained above, you also need to install additional \latex sources
to  benefit from advanced features of \advi.
The command:
\begin{verbatim}
        advi-latex-files --install
\end{verbatim}
will tell you what to do. In particular, you may install \latex source file
at the default path in your \latex installation hierarchy:
\begin{verbatim}
        advi-latex-files --install default
\end{verbatim}
Alternatively, you may specify explicitly where to install those files:
\begin{verbatim}
        advi-latex-files --install <DESTINATION>
\end{verbatim}
The \verb"<DESTINATION>" path should then be included in to your \verb"TEXINPUTS" 
environment variable.

You may also install these additional \latex source files manually. 
The command
\begin{verbatim}
        advi-latex-files --path
\end{verbatim}
tells you where to find those files 
(and you may actually add this path  in your \verb"TEXINPUTS" environment
variable), while the the command: 
\begin{verbatim}
        advi-latex-files --list
\end{verbatim}
will them the files that will be installed.

Finally, you may run
\begin{verbatim}
        advi-latex-files --uninstall
\end{verbatim}
to deinstall the \latex source files, but only if you installed them at the
default location.

\subsection {Rebuilding the documentation}

The documentation, available at \url {http://advi.inria.fr/manual.html}, is
not installed by default.

You may rebuild the document from the sources:
\begin{verbatim}
        make doc.manual
\end{verbatim}
and install it in you \verb"opam" hierarchy:
\begin{verbatim}
        make install.manual
\end{verbatim}

Building the documnentation requires several \latex packages to be installed
\begin{verbatim}
        eepic
        pstricks
        graphics
        babel
        graphicx
        hyperref
        makeidx
        manual
        tabularx
\end{verbatim}
as well as some extra tools: \verb"latexmk", \verb"pandoc", and \verb"hevea"
(see \url {http://hevea.inria.fr/}, also available under opam).

\subsection {Tests and examples}

After successful installation of {\ActiveDVI}, you may have a
look at example presentations provided with the distribution, in the
directories:

\begin{itemize}
%
 \item {\tt test}: a lot of demonstration presentations. From the main
{\ActiveDVI} directory, just type {\tt {cd test; make all}}, and
launch {\tt advi} on files {\tt *.dvi}.
%
 \item {\tt examples}: examples to get inspiration from to easily
write your own talks. A full range of examples, from the simplest one
to the most involved.
%
  \begin{itemize}
%
    \item {\tt basics}: very simple presentations (slightly longer
presentation in sub-directory {\tt cash}).
%
    \item {\tt slitex}: presentations using the {\tt advi-slides}
package. Each presentation has its own directory with a {\tt Makefile}
to automate the process of keeping the presentation up-to-date with
the sources. Nothing difficult here.
%
    \item {\tt seminar}: one presentation using the popular {\tt
seminar} {\latex} package in the {\tt clock} sub-directory.
%
    \item {\tt prosper}: two example presentations using the {\tt
prosper} {\latex} package in sub-directories {\tt LL} and {\tt
Join}. Quite impressive and involved.
%
  \end{itemize}
\end{itemize}


\section{{\ActiveDVI} for the impatient}

\begin{itemize}
%
\item As a previewer, {\ActiveDVI} can preview any correct DVI file.
%
\item As a presenter, {\ActiveDVI} provides some {\latex} packages to
facilitate animations and interaction with the presenter from within
your {\latex} source text. The {\tt advi-slides.sty} package is
designed to be a simple way to build a presentation for {\ActiveDVI}.
%
\end{itemize}

See {\tt examples/slitex/simplistic/} for a simplistic talk example.  There,
the command {\tt \char92pause} is used to make the presenter to stop while
displaying the document. More involved examples can be found in the
directory {\tt examples} of the distribution.

\section{Safety concerns when using the {\ActiveDVI} previewer}

\paragraph{\bf Warning!}

{\ActiveDVI} may execute programs and commands embedded into the DVI
file.  Hence, when playing a DVI file from an untrusted source, you
should run \doctt{advi} with the \doctt{-safer} option that inhibits
the execution of embedded applications. This warning applies in
particular if you choose {\ActiveDVI} as your default {\tt meta-mail}
previewer for the {\tt application/x-dvi} mime-type.

\bigskip

The default safety option is the \doctt{-ask} option: it
tells {\ActiveDVI} to ask the user each time it must launch an
application. (Note that in such a case {\ActiveDVI} asks only once to
launch a given application: it remembers your previous decisions
concerning the command and acts accordingly for the rest of the
presentation.)

The second safety option is the above mentioned \doctt{-safer} option:
it completely inhibits the execution of embedded applications.

The last safety option is \doctt{-exec}: if you call \doctt{advi}
\doctt{-exec}, \doctt{advi} automatically and silently launches all
embedded applications (this is useful to play your own presentations without
the burden of answering yes to {\ActiveDVI}'s questions).

As mentioned, the safe \doctt{-ask} option is the default,
automatically set when nothing has been explicitly specified by the
user. If desired, the default safety option can be set via {\em
initialization files}, either on a system large scale by the machine
administrator (in the file \doctt{/etc/advirc}), on a local scale by
individual users (setting the default policy for that user), or even
on a per directory basis (setting the default policy to show DVI files
in this directory)! (This last option is convenient to gracefully run
your own talks, while still being cautious when running talks from
others.)

\section{Initialization files for {\ActiveDVI}}

\subsection{Syntax of initialization files}

An initialization file for {\ActiveDVI} is simply a text file that
contains options exactly similar to those you can give on the command
line (with the exception of comments, made of a sharp sign
(\doctt{\char35}) followed by some text that is ignored until the end
of line). For instance:
%
\begin{verbatim}
-exec # I know what I mean!
-bgcolor grey16
-fgcolor grey95
\end{verbatim}
is a valid initialization file that sets the safety policy to
\doctt{-exec}, then sets the background and foreground colors to
obtain a nice reverse video effect.

\subsection{Loading initialization files}

Before parsing options on the command line, {\ActiveDVI} loads, in the
order listed below, the following initialization files (nothing
happens if any of them does not exist):

\begin{itemize}
\item system wide initialization file: \doctt{/etc/advirc},
\item user specific initialization files: \doctt{\char126/.advirc} then
\doctt{\char126/.advi/advirc},
\item directory specific initialization file: \doctt{./.advirc}.
\end{itemize}

\subsection{Automatic setting of options}

In addition, the user may load an arbitrary file containing options by
specifying the file path via the command line argument
\doctt{-options-file}.  Hence, \doctt{-options-file} {\tt filename}
loads {\tt filename} when parsing this option to set up the options
contained in {\tt filename} (thus overriding the options set before by
the default \doctt{\char126/.advirc}, \doctt{\char126/.advi/advirc}, or
\doctt{./.advirc}, initialization files).

{\softtarget{otherpage}{%
 \section{Using the {\ActiveDVI} presenter}}}

\subsection{Launching {\ActiveDVI}}

{\ActiveDVI} is invoked with the following command syntax
\begin{verbatim}
advi [options] dvifile [dvifile]
\end{verbatim}

Once {\ActiveDVI} is launched, just press {\bf \verb"?"} to get help
on the keys you can use to control the presenter (type \verb"^f"
(Control-F) to get full screen, \verb"<" or \verb">" to change
the magnification of the text).

When two file names are provided {\ActiveDVI} displays them both: see
section~\ref{multiplex} for details on the use of a secondary DVI file.

\subsection{Command line options}
%
\noindent 
The \doctt{advi} commands recognized the following options:
%
\subsubsection*{Help and info options}
%
\Medskip\noindent\begin{tabular}{ll}
\verb"-v"         & Prints the \doctt{advi} current version and exits \\
\verb"--version"  & Prints the full \doctt{advi} current version and
exits \\
\verb"-help"      & Short command line options help
\end{tabular}
\subsubsection*{Window and display specifications}
%
\Medskip\noindent\begin{tabular}{ll}
\verb"-geometry geom"      & Geometry of \ActiveDVI's window specification
\end{tabular}
%
Geometry \verb"geom" is specified in pixels, using the standard format
for X-Window geometry specifications (i.e:
\verb"WIDTHxHEIGHT[+XOFFSET+YOFFSET]").
\Medskip\noindent\begin{tabular}{ll}
\verb"-fullwidth"          & Adjust the size of the window to the width of the screen
\end{tabular}
\Medskip\noindent\begin{tabular}{ll}
\verb"-nomargins"          & Cancel horizontal and vertical margins \\
\verb"-hmargin dimen"      & Horizontal margin specification (default: 1cm) \\
\verb"-vmargin dimen"      & Vertical margin specification   (default: 1cm) \\
\end{tabular}
Dimensions are specified as numbers optionally followed by two letters
representing units. When no units are given, dimensions are treated
as numbers of pixels. Currently supported units are the standard TeX
units as specified in the TeX book (D. Knuth, Addison-Wesley, (C)
1986):
 `pt' (point), `pc' (pica), `in' (inch), `bp' (big point),
 `cm' (centimeter), `mm' (millimeter), `dd' (didot point),
 `cc' (cicero) and `sp' (scaled point).
Note that dimensions are specified w.r.t the original TeX document,
and do not correspond to what is actually shown on the screen, which
can be displayed at a different resolution than specified in the
original TeX source.
\Medskip\noindent\begin{tabular}{ll}
\verb"-crop"               & Crop the window to the best size (default) \\
\verb"-nocrop"             & Disable cropping
\end{tabular}
%
\subsubsection*{Color specifications}
\Medskip\noindent\begin{tabular}{ll}
\verb"-fgcolor <color>" & Specify the color of the foreground color \\
\verb"-bgcolor <color>" & Specify the color of the background color \\
\verb"-rv" & Specify that reverse video should be simulated by exchanging \\
& the background and foreground colors \\
\verb"-gamma <float>" & Specify gamma correction ($>0.0$) of glyphs\\
\end{tabular}
%
\subsubsection*{Helpers specification}
%
\Medskip\noindent\begin{tabular}{ll}
\verb"-pager" & Specify the name of the pager to launch on a \verb"txt" link\\
\verb"-browser" & Specify the name of the browser to a \verb"html" link\\
\end{tabular}

\subsubsection*{Debugging options}
%
\Medskip\noindent\begin{tabular}{ll}
\verb"--debug" & General debug option \\
\verb"--debug_pages" & Debug page motion \\
\verb"--show_ps" & Print a copy of Postscript sent to gs to stdout \\
\verb"--verbose_image_access" & Change the cursor while loading images\\
\end{tabular}
%
\subsubsection*{Rendering options}
%
\Medskip\noindent\begin{tabular}{ll}
\verb"-A" & Toggle Postscript anti-aliasing \\
\verb"-passive" & Inhibit effects that are visible when redrawing the page
\\ & (Transitions, delays, embedded applications) \\
\end{tabular}
%
\subsubsection*{Safety options}
\Medskip\noindent\begin{tabular}{ll}
\verb"-exec" & Set safety policy to ``always execute embedded applications'' \\
\verb"-ask" & Set safety policy to ``ask user before execution of embedded applications'' \\
\verb"-safer" & Set safety policy to ``never execute embedded applications''
\end{tabular}
%
\subsubsection*{Option files option}
\Medskip\noindent\begin{tabular}{ll}
\verb"-option-file <filename>" & Load {\tt filename} as a file containing a list of options \\
& as given on the command line to \doctt{advi}. \\
\end{tabular}
%
\subsubsection*{Miscelleanous options}
\Medskip\noindent\begin{tabular}{ll}
\doctt{-autoswitch} & Set the {\tt autoswitch} flag, which allows
implicit switch to master \\& on \texttt{usr1} signal (default is off).\\
\end{tabular}

\subsection{Cut and paste}

Text can be copied from the {\ActiveDVI} previewer to another application.
However, this uses the \doctt{XBuffer} and not the \doctt{XSelection}
mechanism.
\begin {itemize}
\item
Shift middle-click copies the current word.
\item
Shift right-click and drag copies the specified region. 
\end {itemize}
Moreover, Shift left-click dump an ASCII representation of the region
under the mouse pointer in the source file. This expects the DVI to
have been instrumented with line numbers of the form
\begin{quote}
\doctt {line:} \docid {line} \docid {file}
\end{quote}
where \docid {line} and \docid {file} are the current source line and
current source file.

The position is exported in ASCII, in the form
\begin{quote}
\verb"#line "\docid{before}, \docid{after} 
\doctt {<<}\docid{prefix}\doctt{>><<}\docid{suffix}\doctt{>>}
\docid{file} 
\end{quote}
where \docid{before} and \docid{after} are the enclosing line numbers, 
\docid{prefix}\docid{suffix} the word constituent surrounding the
mouse position, and file is the name of the current file. 

Line numbers default to $0$ when not found. Note that line numbers may be
inconsistent if there \docdef\special-line commands have not been inserted
close enough to the position.

\subsection{Hyper-text references}

{\ActiveDVI} supports the {\latex} {\tt hyperref} package with both
internal and cross-file references. For cross-file references, it
launches a new \doctt{advi} process to view the target.

{\ActiveDVI} improves the treatment of hyper-refs over conventional previewers,
by emphasizing the hyper-target text of an hyper-link. Thus, an hyper-target
definition:

\docdef\hypertarget \docarg{tag}\docarg{text}

\noindent
should make the activation of the link \docid{text} not only move to
the page where \docid{tag} occurs, but also emphasize the destination
target \docid{tag}.
%
However, since {\docdef\hypertarget} does include its second argument
within the target, we use the following command instead:
\begin{verbatim}
\edef\hyper@quote{\string"}
\edef\hyper@sharp{\string#}

\newcommand{\softtarget}[2]%
  {\special{html:<a name=\hyper@quote#1\hyper@quote>}#2%
   \special{html:</a>}}
\end{verbatim}

(If you are viewing this document with {\softtarget{targethere}{\ActiveDVI}},
you may move over \softlink{targethere}{this area} 
or click on \softlink{otherpage}{this one} to see the effect.)

Similarly, to define a link target we use:
\begin{verbatim}
\newcommand{\softlink}[2]%
  {\special{html:<a href=\hyper@quote\hyper@sharp#1\hyper@quote>}#2%
   \special{html:</a>}}
\end{verbatim}

\subsection{Floating table of contents and thumbnails}

There are two ways to include a floating table of contents while previewing. 

\begin {itemize}

\item {\ActiveDVI} recognizes the reserved hyper-targets
\doctt{advitoc.first} and \doctt{advitoc.last} as markers for the first and
last pages of the table of contents. These pages then become floating,
{\em i.e.} accessible from anywhere in the document with the default
key binding \verb't'. The first stroke on \verb't' shows the first page
of the table of contents. Successive stokes will show the following
pages. (As usual, prefix integer argument may be used to directly
access a specific page of the table of contents.)

The package \softlink{advistyle}{\doctt{advi}} described below
redefines the macro {\docdef\tableofcontents} so that it automatically
inserts the reserved hyper-targets markers around the table of
contents. It also provides two new macros, {\docdef\advitoc} and
{\docdef\endadvitoc}, that serve to insert these markers when the table
of contents is hand-made.

\item If no table of contents markers are found, then {\ActiveDVI}
will compute thumbnails, {\em i.e.} will show the whole set of pages
of the presentation, each page drawn at a smaller scale and packed
with the others on a single page. {\ActiveDVI} computes the scale so
that all the thumbnails fit on one page only, provided that the scale
is less or equal to a maximal value; otherwise, the maximal value
scale is selected and the thumbnail pages spread on several pages. The
default maximal scale value is 5, so that 25 thumbnails can fit on the
same page. This value can be changed using the option
\doctt{-thumbnail\_scale}.

Normally, thumbnails are drawn for all the pages. However, thumbnail
pages can also be defined manually, with an hyper-target whose anchor
is of the form \doctt{/page.\docid{suffix}}. In this case, all the
desired thumbnails must be explicitly marked.

\end{itemize}

By default, the binding \doctt{T} processes thumbnails and the binding
\doctt{t} displays thumbnails if already processed, or shows the table
of contents if available. Otherwise pressing \doctt{t} has no
action. Thumbnails computation is explicit, so that incidentally
hitting the \doctt{t} key does not lead to an unexpected computation,
hence an unexpected delay.

\subsection{Moving around}

See the \softlink{activedvikeys}{key bindings} in the appendix.

\subsection{Scratching on slides}

During the show you can annotate your slides, entering the {\em
scratching} mode. There are two modes, one for writing characters
(the {\em writing mode}, entered by pressing {\tt s} during the show)
and the other to draw lines or figures on the slides (the {\em drawing
mode}, entered by pressing {\tt S}). In each of this modes, you can
enter the {\em scratch setting mode} to set various properties of the
scratching process. See the relevant key bindings for
\softlink{writingkeys}{writing mode} and \softlink{drawingkeys}{drawing
mode} 
in the appendix.

\subsection{Using the laser pointer}

If you press \verb"^X-l" (Control-X then l) the laser pointer appears
on the slide; the pointer sticks to the mouse pointer and allows easy
pointing to parts of the presentation. The laser pointer size and
color can be set on the command line (options {\tt
-laser-pointer-color} and {\tt laser-pointer-size}).

\subsection{Saving slides}

You can save a snapshot of the current slide at any time by pressing
\verb"^X-^S" (Control-X then Control-S). An image file is written
(by default a {\tt png} file). The name of the file produced can be
set via the command line (see {\tt advi -help} for details) or
directly from within the {\latex} source file with commands
\verb"\advisavepageimage" and \verb"\advisavepageimagefile{filename}").

\subsection{Creating events from the {\latex} source file}

{\ActiveDVI} provides the command \verb"\advipushkeys" that provides
key presses to the presenter as if you had pressed it when viewing the
presentation. For instance:
\begin{verbatim}
\advipushkeys{"q"}
\end{verbatim}
ends the talk immediately.

Note that control keys must be encoded inside key strings passed to
{\ActiveDVI}: we use the Emacs textual convention. For instance, the
character ``Control-A'' (ASCII 1) is denoted by the two charcters
\verb"^X" (i.e. a carret character immediately followed by an X).
Hence, the command
\begin{verbatim}
\advipushkeys{"^X^F"}
\end{verbatim}
switches to full screen mode.

\subsection{Using and making special effects}

Presentation examples can be found in the \verb"examples"
directory. Don't miss to play them! Then, feel free to read their
source code and copy the effects they provide.

{\ActiveDVI} can be used as is, but will shine when driven by a user
with a bias towards programming: special effects can easily be realized
by using the {\latex} packages provided with the distribution.

Creative advanced users may program the presenter at various levels, either
using or defining simple {\latex} macros, writing new {\latex} package
files, or by implementing extensions to the previewer itself.

% If you ever succeed at adding a nice special effect that could be of
% a general interest, or if you can provide a nice and fancy presentation
% of your own that can serve as an example to the \ActiveDVI's users
% community, please don't hesitate to drop a mail to the implementors
% at \verb"advi@inria.fr".

\subsection{Viewing multiple files simultaneously}

\label{multiplex}
{\ActiveDVI} can be invoked with several DVI files (currently only two).
The first file is always used as the master file and others are client
files. The user can switch between files explicitly (see 
\softlink{activedvikeys}{key bindings}) or implicitly.
There is an implicit switch from the master to the client file $c$ when an
hyperlink that is not found in the master file can be found in the client file
$c$; there is also a switch from the client $c$ to the master when using the
history stack and the previous event on the stack was an implicit switch
from the master to the client $c$. 

If \doctt{autoswitch} flag is set, there is also an implicit switch to
the master, whenever {\ActiveDVI} receives signal \doctt{usr1} (to mean 
immediate refresh). 

\softtarget{advistyle}{%
 \section{The {\doctt{advi.sty}} {\latex} package}}

{\ActiveDVI} provides some {\latex} packages to facilitate animations and
interaction with the presenter from within your {\latex} source text.

The \doctt{advi.sty} package is the main package to include when writing
a presentation for {\ActiveDVI}. It defines the main set of interactive
commands for {\ActiveDVI} to animate the show. However, there is no
need to load the package if no {\ActiveDVI} special effects are
required for the presentation.

\paragraph{Warning!} Most commands of \doctt{advi.sty} use the {\TeX}
\verb$\special$ command to insert into the DVI output file the
{\ActiveDVI} specific commands that implement their
semantics. Those commands are interpreted by {\ActiveDVI} afterwhile during
the DVI file previewing. Note that a \verb$\special{bla bla}$ command
is equivalent to a \verb$\hbox{}$ for {\TeX}'s mouth, hence it may alter
the document layout accordingly. Thus, be aware that most commands of
the \doctt{advi.sty} package are equivalent to a \verb$\hbox{}$
command as far as the document layout is concerned.

\subsection{Printing the presentation}

The \doctt{advi.sty} package recognizes the special option \doctt{ignore},
which helps the production of a printable version of the
presentation: the \doctt{ignore} option makes the package not to
produce {\ActiveDVI} specials, so that the show can be previewed by
other DVI previewers or turned into Postscript using {\tt dvips}. Of
course, this option disables most effects that cannot be printed,
although some of them are still approximated.

If the \doctt{ignore} option is not set globally, it can be set
locally with the commands {\docdef \adviignore}. However, this will not
prevent all effects, since some decisions are taken when the package is
loaded.

The package also defines the conditional {\docdef \ifadvi} which
evaluates its first argument if {\doctt advi} is not in ignore mode and its
second argument otherwise.

\subsection{Pauses}

{\ActiveDVI} provides partial display of pages (slide
``strip-tease''): the {\ActiveDVI}'s rendering engine stops before the
display of the current page is complete. The corresponding state is
named a {\em pause}. Upon reaching a pause, {\ActiveDVI} may wait for
a specified delay, or for user input.
%
\Medskip\noindent
\docdef \adviwait\docopt{seconds}
\begin{quote}
Wait for \docid{seconds}. If no argument is provided, waits until the
user requests to continue (hitting a key to move to next pause or to
change page).
\end{quote}

\subsection{{\ActiveDVI} records}

{\ActiveDVI} allows (almost) any piece of {\latex} code to be
recorded, and the corresponding DVI code to be rendered later upon
request. To be able to render the code in any order we choose, we must
bind the recorded {\latex} code to a name (called a {\em tag}): this
{\latex} code together with its tag defines the notion of an {\em
{\ActiveDVI} record}.

\paragraph{Warning!} The entire DVI image of an {\ActiveDVI} record
{\em must} fit on a single DVI page. The corresponding check is left
to the writer of the document.

The command defining an {\ActiveDVI} record is as follows:
%
\Medskip\noindent
\docdef \advirecord \doctt{[play]}\docarg{tag}\docarg{latex code}
\\
\docenv{advirecording}{\doctt{[play]}\docarg{tag}\docarg{latex code}%
\docid {text}}

\begin{quote}
This command processes \docid{latex code} and records the
corresponding DVI output, then binds it to the tag \docid{tag}. While
recording, the DVI output is not displayed, unless the option
\doctt{play} is set.

{\ActiveDVI} records may be nested. In this case, the inner record is
bound to its own tag as usual; in addition, if the inner record is
defined with the \doctt{play} option, it is also recorded as a part of
the outer tag record.

If the environment syntax form of {\ActiveDVI} record definition is
used, the \docid{latex code} may contain fragile commands.
\end{quote}

To play an {\ActiveDVI} record, the corresponding DVI must have been
recorded on the current DVI page and before issuing the play command.
With the proviso that these requirements are satisfied, the syntax of
the command to display an {\ActiveDVI} record is as follows:
%
\Medskip\noindent
\docdef \adviplay\docopt{color}\docarg{tag}
\begin{quote}
This command plays the DVI code previously recorded and bound to \docid{tag}.

The optional argument changes the text color to \docid{color} during replay. 
\end{quote}

\subsection{{\ActiveDVI} anchors}

{\ActiveDVI} gives the ability to define {\em {\ActiveDVI}
anchors}: an anchor is specified by (1) an {\ActiveDVI} record, (2) a {\latex}
piece of code that defines the area of the page where the anchor is
active, and (3) an activation method.

The anchor is {\em activated} when some mouse events specified by the
activation method occur in the area. The anchor is {\em de-activated}
when the event that triggered the activation does not hold any more.

The {\ActiveDVI} record associated with the anchor is automatically
rendered \emph{anew} each time the anchor is activated.  The page is
reset to its original appearance when the anchor is de-activated.

The activation method of an anchor may be either \doctt{over} or
\doctt{click}. If the activation method is \doctt{over}, the anchor is
activated whenever the mouse pointer is inside the anchor area;
conversely, the anchor is de-activated when the mouse leaves the
anchor area. If the activation method is \doctt{click}, the anchor is
activated whenever the button is pressed inside the anchor area;
conversely, the anchor is de-activated when the button is released.
%
\Medskip\noindent
\docdef \advianchor\docopt{activation}\docarg{tag}\docarg{text}
\\
\docenv{advianchoring}{\docopt{activation}\docarg{tag}%
\docid {text}}

\begin{quote}
This command defines the DVI rendering of \docarg{text} as an
{\ActiveDVI} anchor area that plays the {\ActiveDVI} record bound to
\docid{tag} when the anchor is activated.

The argument \docid{activation} specifies the activation method of the
anchor.

If the environment syntax form is used, \docid{text} may contain
fragile commands.
\end{quote}

\subsection{Images}

Images can be encapsulated into the presentation using the Caml library
CamlImages provided with the distribution of {\ActiveDVI} (see
section~\ref {advi-graphicx}).

Images can be {\em lighten} by specifying an alpha value (a
floating point number between 0 and 1) that measures the mixing
between the background and the image.

Images can also be {\em blended,} meaning that you can choose the
algorithm that superimposes the image to the background.
Blending modes are reminiscent of the Ghostscript blending options: the blend
mode must be one of the following: {\tt normal}, {\tt multiply},
{\tt screen}, {\tt overlay}, {\tt dodge}, {\tt burn}, {\tt darken},
{\tt lighten}, {\tt difference}, {\tt exclusion}, (\doctt{none} means
unset).

\begin{verbatim}
{\setblend{burn}
  {\setalpha{0.5}
    {\includegraphics[width = 0.7\textwidth]{bar.eps}}}}
\end{verbatim}

{\setblend{burn}
  {\setalpha{0.5}
    {\includegraphics[width = 0.7\textwidth]{bar.eps}}}}

%% {\setblend{normal}
%%   {\setalpha{0.5}
%%     {\includegraphics[width = 0.7\textwidth]{bar.eps}}}}

\subsection{Colors}

\subsubsection*{The \doctt{color} {\latex} package}

In {\ActiveDVI}, colors can be specified with the conventions of the
{\latex} package \doctt{color.sty}, that is, it should either be a
previously defined color or a specification of the form
\docopt{model} \docarg{model color specification}.

For example, the following specifications are all correct:
\begin{quote}
\begin{verbatim}
\color{blue}
\color[named]{Yellow}
\color[rgb]{0.7,0.3,0.8}
\end{verbatim}
\end{quote}

\subsubsection*{Named colors}

Colors can be named using the keyword \docid{named}.  If you use named
colors, the color names are case sensitive and should generally be
capitalized; for instance: \verb"\color[named]{White}" specifies the
white color. Hence, {\tt {\char92}color[named]\{Red\}\{{\it some
text}\}} writes {\it some text} in {\color[named]{Red}{red}}.

The names of available colors can be found in the {\tt dvipsnam.def}
file, generally at location\\
{\tt /usr/share/texmf/tex/latex/graphics/dvipsnam.def}.

To give an idea, the names and colors available on a standard installation of
{\latex} are:

{\color[named]{GreenYellow}{GreenYellow}} $\quad$
{\color[named]{Yellow}{Yellow}} $\quad$
{\color[named]{Goldenrod}{Goldenrod}} $\quad$
{\color[named]{Dandelion}{Dandelion}} $\quad$
{\color[named]{Apricot}{Apricot}} $\quad$

{\color[named]{Peach}{Peach}} $\quad$
{\color[named]{Melon}{Melon}} $\quad$
{\color[named]{YellowOrange}{YellowOrange}} $\quad$
{\color[named]{Orange}{Orange}} $\quad$
{\color[named]{BurntOrange}{BurntOrange}} $\quad$

{\color[named]{Bittersweet}{Bittersweet}} $\quad$
{\color[named]{RedOrange}{RedOrange}} $\quad$
{\color[named]{Mahogany}{Mahogany}} $\quad$
{\color[named]{Maroon}{Maroon}} $\quad$
{\color[named]{BrickRed}{BrickRed}} $\quad$

{\color[named]{Red}{Red}} $\quad$
{\color[named]{OrangeRed}{OrangeRed}} $\quad$
{\color[named]{RubineRed}{RubineRed}} $\quad$
{\color[named]{WildStrawberry}{WildStrawberry}} $\quad$
{\color[named]{Salmon}{Salmon}} $\quad$

{\color[named]{CarnationPink}{CarnationPink}} $\quad$
{\color[named]{Magenta}{Magenta}} $\quad$
{\color[named]{VioletRed}{VioletRed}} $\quad$
{\color[named]{Rhodamine}{Rhodamine}} $\quad$
{\color[named]{Mulberry}{Mulberry}} $\quad$

{\color[named]{RedViolet}{RedViolet}} $\quad$
{\color[named]{Fuchsia}{Fuchsia}} $\quad$
{\color[named]{Lavender}{Lavender}} $\quad$
{\color[named]{Thistle}{Thistle}} $\quad$
{\color[named]{Orchid}{Orchid}} $\quad$
{\color[named]{DarkOrchid}{DarkOrchid}} $\quad$

{\color[named]{Purple}{Purple}} $\quad$
{\color[named]{Plum}{Plum}} $\quad$
{\color[named]{Violet}{Violet}} $\quad$
{\color[named]{RoyalPurple}{RoyalPurple}} $\quad$
{\color[named]{BlueViolet}{BlueViolet}} $\quad$

{\color[named]{Periwinkle}{Periwinkle}} $\quad$
{\color[named]{CadetBlue}{CadetBlue}} $\quad$
{\color[named]{CornflowerBlue}{CornflowerBlue}} $\quad$
{\color[named]{MidnightBlue}{MidnightBlue}} $\quad$

{\color[named]{NavyBlue}{NavyBlue}} $\quad$
{\color[named]{RoyalBlue}{RoyalBlue}} $\quad$
{\color[named]{Blue}{Blue}} $\quad$
{\color[named]{Cerulean}{Cerulean}} $\quad$
{\color[named]{Cyan}{Cyan}} $\quad$
{\color[named]{ProcessBlue}{ProcessBlue}} $\quad$

{\color[named]{SkyBlue}{SkyBlue}} $\quad$
{\color[named]{Turquoise}{Turquoise}} $\quad$
{\color[named]{TealBlue}{TealBlue}} $\quad$
{\color[named]{Aquamarine}{Aquamarine}} $\quad$
{\color[named]{BlueGreen}{BlueGreen}} $\quad$

{\color[named]{Emerald}{Emerald}} $\quad$
{\color[named]{JungleGreen}{JungleGreen}} $\quad$
{\color[named]{SeaGreen}{SeaGreen}} $\quad$
{\color[named]{Green}{Green}} $\quad$
{\color[named]{ForestGreen}{ForestGreen}} $\quad$

{\color[named]{PineGreen}{PineGreen}} $\quad$
{\color[named]{LimeGreen}{LimeGreen}} $\quad$
{\color[named]{YellowGreen}{YellowGreen}} $\quad$
{\color[named]{SpringGreen}{SpringGreen}} $\quad$

{\color[named]{OliveGreen}{OliveGreen}} $\quad$
{\color[named]{RawSienna}{RawSienna}} $\quad$
{\color[named]{Sepia}{Sepia}} $\quad$
{\color[named]{Brown}{Brown}} $\quad$
{\color[named]{Tan}{Tan}} $\quad$

{\color[named]{Gray}{Gray}} $\quad$
{\color[named]{Black}{Black}} $\quad$

\subsubsection*{The CMYK specifications of colors}

You may also explicitly use a CMYK (Cyan, Magenta, Yellow, Black)
specification. In this case the cyan, magenta, yellow and black values
follow the \docid{cmyk} keyword, and are given as a list of four
integers in the range 0.0 .. 1.0.
For instance, {\tt {\char92}color[cmyk]\{0,1,0,0\}} is a valid specification
for {\color[cmyk]{0,1,0,0}{magenta}}.

\subsubsection*{The RGB specifications of colors}

RGB (Red, Green, Blue) specifications are similar to the CMYK
specifications: following the \docid{rbg} keyword, the red, green, or
blue color values, are given as floating point numbers in the range
0.0 .. 1.0.  Hence, {\tt{{\char92}color[rgb]\{1.0,0.0,0.0\}}} is a
valid specification for {\color[rgb]{1.0,0.0,0.0}{red}}.

\subsubsection*{The X -Window System colors}

{\ActiveDVI} provides the package \doctt{xwindows-colors}, an
extension to the \doctt{color} package, that defines a large set of
the X Window System color names, as found in the file {\tt rgb.txt}
of a typical X installation (this file is generally located on {\tt
/usr/X11R6/lib/X11/rgb.txt}). To know which colors are available look
at the source file of the package {\tt xwindows-colors.sty} in the
directory {\tt tex} of the distribution.

\subsection{Background}

You can modify the background of your presentation in the {\latex}
source of the pages. Background can be defined either as a plain
color, as an image, or as a gradient (or as a combination of these!).

You can specify a global option to the background settings, so that these
settings are used for the remaining pages of the presentation
(otherwise the presenter resets the background options at each new page).

To modify the background of your presentation, you can:

\begin{itemize}
  \item define the background color,
  \item define a gradient function to be run on the background (or a
   defined area of the background) using the predefined color gradients,
  \item add a background image (which can be alpha-blended on top of
   the background color).
\end{itemize}

If these options are used \emph{together}, they are applied in this
order: first the solid background color is drawn, then the gradient
function is applied, finally the image is drawn on the resulting
background.
%
\Medskip\noindent
\docdef \advibg \doctt {[global]}\docarg{decl}
\begin{quote}
where \docid{decl} is a list of settings of the following from:

\dockey[none]{color}{color}
\begin{quote}
Set the background color to \docid{color}.
If \docid{color} is \doctt{none} this unsets the background color. 
Otherwise, \docid{color} must follow the notation above to designate colors.
\end{quote}

\dockey[none]{image}{file}
\begin{quote}
Use the image found in \docid{file} as background (\doctt{none} means unset).
\end{quote}

\dockey[auto]{fit}{fit style}
\begin{quote}
Fit the background image according to \docid{style}, which may be 
one of the following keywords:
\begin{center}
\tt auto \hfill {\rm or} \hfill
\begin{tabular}{l@{\qquad}c@{\qquad}r}
topleft    & top     & topright\\
left       & center  & right \\
bottomleft & bottom  & bottomright\\
\end{tabular}
\end{center}
The {\tt auto} fit style means scaling the image as desired in both directions
so that it fits the entire page. Other styles only force the same scaling
factor in both directions:

\begin {itemize}

\item Corner-styles means set the image in the corresponding corner
and scale it to cover the entire page.

\item {\tt center} means set the image in the center of the page and
scale it to cover the entire page.

\item Segment-styles means adjust the image and the page on the
segment (in which case, the image may not completely cover the page on
the opposite side).

\end {itemize}
\end{quote}

\dockey[none]{alpha}{float}
\begin{quote}
Set the alpha channel factor for the background image to \docid{float}
(\doctt{none} means unset). An alpha factor of $0$ means that the
image is not visible at all; conversely, an alpha factor of $1$ means that the
image covers the background.
\end{quote}

\dockey[none]{blend}{blend mode}
\begin{quote}
Set the blend mode to \docid{blend mode}, which are reminiscent of
Ghostscript blending options. The blend mode should be one of the
following: {\tt normal}, {\tt multiply}, {\tt screen}, {\tt overlay},
{\tt dodge}, {\tt burn}, {\tt darken}, {\tt lighten}, {\tt
difference}, {\tt exclusion}, (\doctt{none} means unset).
\end{quote}

\dockey[none]{gradient}{function}
\begin{quote}
Set the gradient function to \docid{function}, one of the
predefined functions that convert the plain background color into a
color gradient from the chosen color \verb"colorstart" to the color
\verb"colorstop" (which is white by default). Available gradients are:
\begin {itemize}

\item \verb|hgradient| horizontal gradient (the gradient is a serie of
vertical lines),
\item \verb|vgradient| vertical gradient (the gradient is a serie of
horizontal lines),
\item \verb|d1gradient| first bissector gradient (the gradient is a serie of
lines which are parallel to the first bissector),
\item \verb|d2gradient| second bissector gradient (the gradient is a serie of
lines which are parallel to the second bissector),
\item \verb|cgradient| centered gradient (the gradient is a serie of
concentric squares),
\item \verb|circgradient| circular gradient (the gradient is a serie
of concentric circles).
\end {itemize}
\end{quote}

\dockey[white]{colorstart}{color}
\begin{quote}
Set the starting color of the gradient. When left unspecified defaults to
white.
\end{quote}

\dockey[background]{colorstop}{color}
\begin{quote}
Set the end color of the gradient. When left unspecified defaults to
the background color.
\end{quote}

\dockey[0]{xstart}{int}
\begin{quote}
Set the abscissa of the lower left point of the area where the
gradient is drawn.
\end{quote}

\dockey[0]{ystart}{int}
\begin{quote}
Set the ordinate of the lower left point of the area where the
gradient is drawn.
\end{quote}

\dockey[1.0]{width}{float}
\begin{quote}
Set the width of the area where the gradient is drawn.
The width is a number in the range \verb|[0 .. 1]| that gives the
ratio of the area width with respect to the page width (hence $0.0$
means a null width and $1.0$ means the entire page width).
\end{quote}

\dockey[1.0]{height}{float}
\begin{quote}
Set the height of the area where the gradient is drawn.
The width is a number in the range \verb|[0 .. 1]| that gives the
ratio of the area height with respect to the page height (hence $0.0$
means a null height and $1.0$ means the entire page height).
\end{quote}

\dockey[0.5]{xcenter}{float}
\begin{quote}
For a centered or circular gradient, set the abscissa of the center
point of the gradient into the gradient area. {\tt xcenter} is a ratio
of the gradient area's width. It defaults to $0.5$, meaning
the middle of the gradient area width.
\end{quote}

\dockey[0.5]{ycenter}{float}
\begin{quote}
For a centered or circular gradient, set the ordinate of the center
point of the gradient into the gradient area. {\tt ycenter} is a ratio
of the gradient area's height. It defaults to $0.5$, meaning
the middle of the gradient area height.
\end{quote}

\goodbreak
\dockey{none}{}
\begin{quote}
Unset all background parameters.
This key must appear on its own, no arguments or keys are allowed. 
\end{quote}
The optional parameter \doctt{global} indicates that
the definition is global and will affect the following pages, 
as well as the current page.

By default, the background settings only affect the current page.
\end{quote}

\subsection{Transitions}

\docdef \advitransition \doctt {[global]}\docarg {decl}
\begin{quote}
where \docid{decl} is a list of settings of the following from:

\dockey{none}{} or \dockey{slide}{} or \dockey{block}{} or \dockey{wipe}{}
\begin{quote}
Set the transition mode to the corresponding key.
One of this key is mandatory (if several are provided the last one is
selected).
\end{quote}

\dockey{from}{direction}
\begin{quote}
Make the transition come from \docid{direction}. Directions should be
one of the following:
\begin{center}
\tt
\begin{tabular}{l@{\qquad}c@{\qquad}r}
topleft    &  top     & topright\\
left       &  center  & right \\
bottomleft &  bottom  & bottomright\\
\end{tabular}
\end{center}
\end{quote}
The default direction, to be used when no local or global direction has been
specified, is determined dynamically: \verb"right" when coming from
previous page, \verb"left" when coming from next page, and \verb"top"
otherwise.

\dockey{steps}{n}
\begin{quote}
Make the transition in \docid{n} steps. 
\end{quote}
As for \docdef \advibg, the optional parameter \doctt{global} indicates that
the definition is global and will affect the following pages, 
as well as the current page. 

By default, the transition definitions affect the current page only. 
\end{quote}
%
\Medskip\noindent
\docdef \advitransbox \docarg {key=val list}\docarg{hbox material}
\begin{quote}
where \docid {key=val list} is as above and \docarg{hbox material} is
whatever can follow an \docdef \hbox{} command. In particular, the material
may contain verbatim commands, since as for the \docdef\hbox{} it is
parsed incrementally.
\end{quote}

The optional parameter \doctt{global} indicates that the definition is
global and will affect the following pages, as well as the current
page.

By default, the transition affects the current page only.

\subsection{Embedded applications}

To animate your show, {\ActiveDVI} can launch arbitrary applications
you need.

\subsubsection{Launching embedded applications}

The {\latex} command to launch an application during the presentation is 
%
\Medskip\noindent
\docdef \adviembed\docopt{key=value list}\docarg{command} 

\begin{quote}
where \docid{key=value list} is a list of bindings of the following kind:

\dockey{name}{name}
\begin{quote}
Allows to refer to the embedded application as \docid{name}. 
Anonymous applications have actually the default name \doctt{anonymous}.
\end{quote}

\dockey{ephemeral}{name}

\begin{quote}
This is the default case: the application is specific to a given
page. An \doctt{ephemeral} application is automatically launched
whenever the page is displayed, and automatically killed when the page
is turned.
\end{quote}

\dockey{persistent}{name}

\begin{quote}
A \doctt{persistent} application is launched only once and keeps
running in the background; however, {\ActiveDVI} automatically
hides and shows the window where the application runs, so that the
application is visible only on the page where it has been launched.
\end{quote}

\dockey{sticky}{name}
\begin{quote}
A \doctt{stiky} application is launched only once, keeps running, and
remains visible when turning pages. It is also resized and moved as
necessary to fit the page size.
\end{quote}

\dockey{raw}{name}
\begin{quote}
A \doctt{raw} application is launched each time its embedding command
is encountered. A \doctt{raw} application is not managed {\em
automatically} by {\ActiveDVI}, except for the initial launching and
the final clean-up that occurs when {\ActiveDVI} exits; hence, you can
completely monitor the \doctt{raw} applications graphical behavior,
using the \doctt{advikillembed} command and the window mapping
facilities for \doctt{raw} applications described below.
\end{quote}

\dockey{width}{dim}\\
\dockey{height}{dim}

\begin{quote}
The application takes \docid{dim} width (respectively height) space in 
{\latex}. Both values default to 0pt.

These dimensions are also substituted for all occurrences of \verb"@g" in
the command string.
\end{quote}

\end{quote}

\subsubsection{Monitoring embedded applications}

To monitor embedded applications, {\ActiveDVI} provides the
\doctt{advikillembed} primitive to send a signal to any named embedded
application. For \doctt{raw} applications, there are additional
functions to map or un-map the window allocated to a named \doctt{raw}
application. Mapping or un-mapping windows of non-\doctt{raw}
applications is unspecified, since it may interfere in a non trivial
way with {\ActiveDVI}'s automatic treatment of those applications.

\subsubsection*{Monitoring a single application}
%
\Medskip\noindent
\docdef \advikillembed \docarg{name}
\begin{quote}
Kill the embedded application named \docid{name}.
An optional signal value or symbolic name can be given to send to the
designed process: for instance, {\tt {\char92}advikillembed[SIGUSR1]\{clock\}}
will send the {\tt SIGUSR1} signal to the embedded application named
{\tt clock}.

Signal value defaults to -9.
\end{quote}
%
\Medskip\noindent
\docdef \advimapembed \docarg{name}
\begin{quote}
Map the window of the (\doctt{raw}) embedded application named \docid{name}.
\end{quote}
%
\Medskip\noindent
\docdef \adviunmapembed \docarg{name}
\begin{quote}
Un-map the window of the (\doctt{raw}) embedded application named \docid{name}.
\end{quote}

\subsubsection*{Monitoring a group of embedded applications}

The primitives \doctt{advikillallembed}, \doctt{advimapallembed}, and
\doctt{adviunmapallembed} behave the same as their non-{\tt all}
counterparts, except that they operate on all the applications that
have been launched with the given name.
%
\Medskip\noindent
\docdef \advikillallembed \docarg{name}
\begin{quote}
Similar to \doctt{advikillembed}
but kill all the embedded applications named \docid{name}.
\end{quote}
%
\Medskip\noindent
\docdef \advimapallembed \docarg{name}
\begin{quote}
Map the windows of all the (\doctt{raw}) embedded applications named
\docid{name}.
\end{quote}
%
\Medskip\noindent
\docdef \adviunmapallembed \docarg{name}
\begin{quote}
Un-map the windows of all the (\doctt{raw}) embedded applications named
\docid{name}.
\end{quote}
%
\subsection{Active anchors}
%
Active anchors are annotated pieces of text that get associated
activation records. To define an active anchor, the command is
%
\Medskip\noindent
\docdef \advianchor \docopt{decl}\docarg{tag}\docarg {text}
\\
\docenv{advianchor}{\docopt{decl}\docarg{tag} \docid {text}}
\begin{quote}
The text is first displayed as usual, then the anchor is drawn
according to the style given by \docid{decl}, and made active. 
Its activation, which depends on the mode given by \docid{decl}, 
will play the record named \docid{tag}. 

The declarations \docid{decl} are of the following form:
\begin{description}
\item [\doctt{over}, \doctt{click}, \textrm{or} \doctt{stick}]\indent\par 
The mode \doctt{stick} plays the tag \docid{tag} on
click. The mode \doctt{click} is similar, except that it restores the
previous state when leaving the anchor area. The mode \doctt{over} is as
\doctt{click} but display the \docid{tag} when the mouse is \doctt{over}
the anchor instead of waiting for a click.
\item [\doctt{box}, \doctt{invisible}, or \doctt{underline}]\indent\par
this defines the style in which the anchor should be drawn.
The default style is \doctt{box}.
\end{description}
In the environment form, \docid{text} may contain fragile commands. 
\end{quote}
%
\Medskip\noindent
\docdef \adviemphasize \docopt{color}\docarg {text}
\begin{quote}
This makes an invisible anchor around \docid{text}, which when activated
will redraw text in a box colored with \docid{color}, which defaults to
\doctt{yellow}. 
\end{quote}

\subsection{Postscript specials}

{\ActiveDVI} can deal with most of PStricks by calling {\tt
ghostscript} on included Postscripts. Basic change of coordinates are
implemented, but this feature remain fragile, as {\ActiveDVI} must in
turn call {\tt ghostscript} to get the new coordinates. Also,
rotations will definitely not work for text, which is rendered by
{\ActiveDVI} and cannot be rotated.

\subsubsection{Overlays}

\index{overlays}

The {\tt overlay} class implements overlays with PStricks. By contrast,
{\ActiveDVI} implements overlays directly, using records and plays.
This is more efficient, and of course more natural.
(In fact, {\ActiveDVI} chooses the cumulative semantics of overlays,
displaying all layers below the current overlay.)

The {\tt xprosper} style, derived from the {\tt prosper} class, uses 
the {\tt overlay} class and works with {\ActiveDVI} in exactly the same
way (relaxing the \verb|\overlay@loop| macro inhibits all layers, but the first
page).

\subsubsection{PStricks known to work}

{\ActiveDVI} can deal with main PStricks. In particular, 
the following work
\begin {itemize}

\item[+]
Simple drawings, such as \verb"\psframe", \verb"\ovalnode", \ldots.

\item[+]
Connections between nodes {\tt \string \ncline}, {\tt \string \ncarc},
also works.

\item[+]
Labels over arrows \verb"\Aput", \verb"Bput", etc.

\item[+] 
{\tt\string\SpecialCoor}, {\em i.e.} commands of the form
{\tt \string \rput (A)\{bla bla\}} works where {\tt A} is a node name. 

\item[+]
Embeddded-Postcript figures, including scaling.

\end {itemize}
Other PStricks may or may not work.


\softtarget{advi-slides}{%
 \section{The {\doctt{advi-slides.sty}} {\latex} package}}

{\ActiveDVI} provides this specialized {\latex} package to facilitate
writing presentation slides in the spirit of SliTeX. See examples in
the directory {\tt examples/slitex}.

\section{Auxiliary {\latex} packages}

\subsection{The \texttt{superpose} package}

This package allows superposition of horizontal material, creating 
the smallest horizontal box that fits all of the superpositions.
\begin{verbatim}
\usepackage{superpose}
\end{verbatim}
The package defines a single environment:
%
\Medskip\noindent
\docenv{superpose}{\docopt{alignment}\docid {list}}
\begin{quote}
The \docid{alignment} can be one the letters \doctt c (default value), \doctt
l, or \doctt r.

Items of the \docid{list} are separated by \verb"\\" as in tabular
environments. 
Each item should be a horizontal material.
\end{quote}

\subsection{The \texttt{bubble} package}

This package draws bubbles over some text. 
\begin{verbatim}
\usepackage{bubble}
\usepackage[ps]{bubble}
\end{verbatim}
By default bubbles are produced using the \doctt{epic} and \doctt{eepic}
packages, for portability. However, for better rendering and easier 
parameterization, bubbles can also be drawn using the \doctt{ pst-node}
package of the PStricks collection. This is what the \doctt{ps} option
is designed for.

The package defines a single command:
%
\Medskip\noindent
\docdef \bubble \docopt{key=value list}\docarg{anchor}%
    \docopt{ps options}\docpar {pos}\docarg{text}
\begin{quote}
The \docid{key=value list} is a list of bindings of the following kind:

\dockey[yellow]{bg}{color}
\begin{quote}
The background color for annotations. 
\end{quote}

\dockey[yellow]{unit}{dim}
\begin{quote}
Set the package unit to \docid{dim}.
\end{quote}

\dockey[c]{col}{colspec}
\begin{quote}
Where \docid{colspec} is a column specification for the tabular
environment. Moreover, the following abbreviations are recognized:

\smallskip
\tt
\begin{tabular}{|l|l|}
\hline
\rm\em key & \rm expands to \\
\hline
c &  col=c\\
l &  col=l\\
r &  col=r\\
p=\docid{w} & col=p\docarg{w}\\
%\cline{1-2}
\hline
\end{tabular}
\begin{tabular}{|l|l|}
\hline
\rm\em key & \rm expands to \\
\hline
C &  col=\{>\{\$\}c<\{\$\}\}\\
L &  col=\{>\{\$\}l<\{\$\}\}\\
R &  col=\{>\{\$\}r<\{\$\}\}\\
P\docid{w} &  \doctt{col=\{>\{\$\}p\docarg{w}<\{\$\}\}}\\
%\cline{1-2}
\hline
\end{tabular}
\end{quote}

\docid{pos} is the optional relative position of the annotation, 
it defaults to $1,1$, and is counted in the package units. 

\docid{ps options} are passed to the command \docdef\psset) 
in \doctt{ps} mode and ignored otherwise. 
\end{quote}
Parameters (color and tabular columns specifications) can also 
be set globally using the command:
%
\Medskip \noindent
\docdef \bubbleset\docarg{key=value list}

\subsection{The \texttt{advi-annot} package}

This package uses active anchors and the bubbles package to 
provide annotations by raising a bubble when the cursor is over the anchor.

The package defines a single command
%
\Medskip\noindent
\docdef \adviannot \docopt{key=value list}\docarg{anchor}%
    \docopt{ps options}\docpar {pos}\docarg{text}
\begin{quote}
whose options are identical to those of the {\docdef \bubble} macro; however
the bubble appears within an active anchor.
\end{quote}

%Annotation location can also be invisible \advi{C}{\tt{ }}instead of visible.

\subsection{The {\doctt{advi-graphicx}} {\latex} package}
\label {advi-graphicx}

This 3-lines long package loads the {\tt graphicx.sty} package and provides
declarations so that JPEG, EPS, TIF, TIFF source images can be embedded:
{\ActiveDVI} will preview these images directly while other drivers
will translate them on demand. 

\appendix

\section{Limitations}

\subsection*{Postscript Fonts}

Postscript fonts are not natively handled by {\ActiveDVI}. You must
use the command {\tt dvicopy} to expand those virtual fonts to base
fonts before visualization with {\ActiveDVI}.
(For instance, {\tt dvicopy talk.dvi talk.expanded; advi talk.expanded} very
often does the trick.)

\subsection*{In-lined Postscript and Ghostscript}

PS relies on {\tt ghostscript} to display Postscript in-lined specials.
However, some earlier releases of {\tt ghostscript} 
implements the Postscript {\tt flushpage} command as a {\tt XFlush} call
which does not force the evaluation of commands, and thus makes the 
synchronization between {\tt ghostscript} and {\ActiveDVI} drawings
uncontrollable. In this case, the interleaving of in-lined postscript 
and other material may be inconsistent. 

Fortunately, recent versions of ghostscript $(> 6.5)$ have
fixed this problem by using {\tt XSync(false)} instead.
If you use those versions of ghostscript, in-lined specials should be
correctly rendered.

Unfortunately, some releases of version 6.5x also carry a small but
fatal bug for \ActiveDVI, that will hopefully be fixed in future
releases. A workaround is available here
\url{http://cristal.inria.fr/~remy/ghostscript/}.

\subsection*{In-lined Postscript change of coordinates}

So far, the implementation of in-lined Postscript does not correctly handle
complex change of coordinates.  (See PStricks section). 

\section{Reporting bugs}

Please, send bug reports to 
\url{mailto:advi@inria.fr}. 
\\
See \url{http://gallium.inria.fr/advi} for up to date information.


\softtarget{activedvikeys}{%
 \section{{Key bindings}}}


\let \Newpage \newpage

\def \newpage {}


{\ActiveDVI} recognizes the keystrokes listed below when typed in its window.
Some keystrokes may optionally be preceded by a number, called \arg
below, whose interpretation is keystroke dependant. If \arg is unset,
its value is 1, unless specified otherwise.

{\ActiveDVI} maintains an history of previously visited pages
organized as a stack.
Additionnally, the history contains marked pages which are stronger than
unmarked pages. 

\Stretch\newpage\Stretch


Advi recognizes the keystrokes listed below when typed in its window.
Some keystrokes may optionally be preceded by a number, called \arg
below, whose interpretation is keystroke dependant. If \arg is unset,
its value is 1, unless specified otherwise.

Advi maintains an history of previously visited pages organized as a stack. 
Additionnally, the history contains marked pages which are stronger than
unmarked pages. 

\vspace{\stretch{1}}

\newpage

\vspace{\stretch{1}}

\subsection*{Survival command kit}

\noindent
\begin{tabularx}{\linewidth}{clcX}
\ikey{?}{info}{This quick info and key bindings help.}
\ikey{q}{quit}{End of show.}
\ikey{space}{continue}
{Move forward (\arg pauses forward if any, or do as \key{return} otherwise).}
\end{tabularx}

\vspace{\stretch{1}}

\subsection*{Moving to pages}

\noindent
\begin{tabularx}{\linewidth}{clcX}
\ikey{n}{next}
{Move \arg physical pages forward, leaving the history unchanged.}
\ikey{p}{previous}
{Move \arg physical pages backward, leaving the history unchanged.}
\ikey{,}{begin}{Move to the first page.}
\ikey{.}{end}{Move to the last page.}
\ikey{g}{go}
{If \arg is unset move to the last page.
 If \arg is the current page do nothing.
 Otherwise, push the current page on the history as marked, and move
 to physical page \arg.}
\end{tabularx}

\vspace{\stretch{1}}
\subsection*{Switching views}

\noindent
\begin{tabularx}{\linewidth}{clcX}
\ikey{w}{switch}{Switch view between master and client (if any)}
\ikey{W}{sync}{Goto page of client view corresponding to page of master view}
\ikey{\char94W}{autoswitch}{Toggle autoswitch flag}
\end{tabularx}

\vspace{\stretch{1}}

\subsection*{Table of contents}

\noindent
\begin{tabularx}{\linewidth}{clcX}
\ikey{T}{Thumbnails}{Process thumbnails.}
\ikey{t}{toc}{Display thumbnails if processed, or floating table of contents
 if available, or do nothing.}
\end{tabularx}

\vspace{\stretch{1}}

\vspace{\stretch{1}}

\newpage

\subsection*{Moving to pauses}

\noindent
\begin{tabularx}{\linewidth}{clcX}
\ikey{N}{next pause}{Move \arg pauses forward (equivalent to continue).}
\ikey{P}{previous pause}{Move \arg pauses backward.}
\end{tabularx}

\vspace{\stretch{1}}

\subsection*{Adjusting the page size}

\noindent
\begin{tabularx}{\linewidth}{clcX}
\ikey{\char94F}{fullscreen}{Adjust the size of the page to fit the
entire screen or reset the page to the default size (this is a toggle).}
\ikey{$<$}{smaller}{Scale down the resolution by scalestep (default
\tiny{$\sqrt{\sqrt{\sqrt 2}}$}).}
\ikey{$>$}{bigger}{Scale up the resolution by scalestep (default
\tiny{$\sqrt{\sqrt{\sqrt 2}}$}).}
\ikey{\char35}{fullpage}{Remove margins around the page and change
the resolution accordingly.}
\ikey{c}{center}{Center the page in the window, and resets the default
resolution.}
\end{tabularx}

\vspace{\stretch{1}}

\subsection*{Redisplay commands}

\noindent
\begin{tabularx}{\linewidth}{clcX}
\ikey{r}{redraw}{Redraw the current page to the current pause.}
\ikey{R}{reload}{Reload the file and redraw the current page.}
\ikey{\char94L}{redisplay}{Redisplay the current page to the first
pause of the page.}
\ikey{a}{active/passive}{toggle advi effects (so that reloading is silent).}
\end{tabularx}

\vspace{\stretch{1}}

\newpage

\vspace{\stretch{1}}

\subsection*{Moving the page in the window}

\noindent
\begin{tabularx}{\linewidth}{clcX}
\ikey{h}{page left}
{Moves one screen width toward the left of the page. Does nothing if the
  left part of the page is already displayed}
\ikey{l}{page right}
{Moves one screen width toward the right of the page. Does nothing if the
  right part of the page is already displayed}
\ikey{j}{page down}
{Moves one screen height toward the bottom of the page. Jumps to the top of
  next page, if there is one, and if the bottom of the page is already
  displayed.}
\ikey{k}{page up}
{Moves one screen height toward the top of the page. Jumps to the bottom
 previous page, if there is one, and if the top of the page is already
 displayed.}
\ikey{\char94 left button}{move page}
{A black line draws the page borders; moving the mouse then moves the
page in the window.}
\end{tabularx}

\vspace{\stretch{1}}

\newpage

\vspace{\stretch{1}}

\subsection*{Using the navigation history stack}

\noindent
\begin{tabularx}{\linewidth}{clcX}
\ikey{return}{forward}
{Push the current page on the history stack, and move forward n physical pages.}
\ikey{tab}{mark and next}
{Push the current page on the history as marked, and move forward n
physical pages.}
\ikey{backspace}{back}
{Move \arg pages backward according to the history. The history stack
is poped, accordingly.}
\ikey{escape}{find mark}
{Move \arg marked pages backward according to the history.
 Do nothing if the history does no contain any marked page.}
\end{tabularx}

\vspace{\stretch{1}}

\subsection*{Writing and drawing on the slide}

\noindent
\begin{tabularx}{\linewidth}{clcX}
\ikey{s}{scratch}{Give a pencil to type characters on the page.}
\ikey{S}{scratch}{Give a spray can to draw on the page.}
\end{tabularx}

\subsection*{Dealing with caches}

\noindent
\begin{tabularx}{\linewidth}{clcX}
\ikey{f}{load fonts}{Load all the fonts used in the document. By
default, fonts are loaded only when needed.}
\ikey{F}{make fonts}
{Does the same as \key{f}, and precomputes the glyphs of all the
characters used in the document.
This takes more time than loading the fonts, but the pages are drawn faster.}
\ikey{C}{clear}{Erase the image cache.}
\end{tabularx}

\vspace{\stretch{1}}







\softtarget{writingkeys}{%
 \section{{Key bindings for scratch writings}}}

\subsection*{Entering scratch writing mode}

Press \key{s} to enter scratch writing; the cursor is modified and
you must click somewhere on the page to start writing text
there. Before clicking, you can
\begin{citemize}
 \item press \key{?} to get help,
 \item press \key{\char94 G} to quit scratching immediately,
 \item press \key{Esc}
 to enter the scratch writing settings mode and tune the font and font size.
\end{citemize}

\subsection*{Survival command kit when scratch writing}

{\ActiveDVI} recognizes the following keystrokes when scratch writing
on the page.

\noindent
\begin{tabularx}{\linewidth}{clcX}
\ikey{\char94 G}{quit}{End of scratch writing.}
\ikey{Esc}{settings}{Enter the scratch writing settings mode.}
\end{tabularx}

In the scratch writing settings mode, the cursor is modified and you
can set some charateristics of the scratch writing facility.
When in doubt, press

\begin{citemize}
 \item press \key{?} to get help,
 \item press \key{\char94 G} to quit scratching immediately,
 \item press \key{Esc} to quit the setting mode.
\end{citemize}

\vspace{\stretch{1}}

\newpage

\subsection*{Scratch writing settings mode keys}

When in the scratch writing settings mode, the following keys have the
following respective meanings:

\noindent
\begin{tabularx}{\linewidth}{clcX}
\ikey{$>$}{greater}{Increments the scratch font size.}
\ikey{$<$}{smaller}{Decrements the scratch font size.}
\ikey{b}{blue}{Set the color of the font to blue.}
\ikey{c}{cyan}{Set the color of the font to cyan.}
\ikey{g}{green}{Set the color of the font to green.}
\ikey{k}{black}{Set the color of the font to black.}
\ikey{m}{magenta}{Set the color of the font to magenta.}
\ikey{r}{red}{Set the color of the font to red.}
\ikey{w}{white}{Set the color of the font to white.}
\ikey{y}{yellow}{Set the color of the font to yellow.}
\ikey{B}{more blue}{Increment the blue component of the color.}
\ikey{G}{more green}{Increment the green component of the current color.}
\ikey{R}{more red}{Increment the red component of the current color.}
\ikey{$+$}{positive increment}{Set the color increment to positive.}
\ikey{$-$}{negative increment}{Set the color increment to negative.}
\ikey{$?$}{help}{Give the list of settings available.}
\ikey{Esc}{quit}{Quit the sratch writing settings mode.}
\end{tabularx}

\subsection*{Setting the scratching font size}

Just press \key{Esc} to enter the scratch writing settings mode, then
\key{$>$} or \key{$<$} to increment or decrement the font size; then
press \key{Esc} again, to leave the scratch writing settings mode and
continue to write on the page with the new font size.

\vspace{\stretch{1}}


\softtarget{drawingkeys}{%
 \section{{Key bindings for scratch drawings}}}

Advi recognizes the keystrokes listed below when drawing on the slide.

\subsection*{Survival command kit when scratch drawing}

\noindent
\begin{tabularx}{\linewidth}{clcX}
\ikey{\char94G}{quit}{End of scratch drawing.}
\ikey{Esc}{settings}{Enter the scratch drawing settings mode.}
\end{tabularx}

In the scratch drawing settings mode, the cursor is modified and you
can set some charateristics of the scratch drawing facility.


\subsection*{Scratch drawing settings mode keys}

When in the scratch drawing settings mode, the following keys have the
following respective meanings:

\subsubsection*{General scratch drawing settings keys}

\begin{itemize}
 \item press \key{?} to get help,
 \item press \key{Esc} to quit the settings mode,
 \item press \key{\char94g} to quit scratching immediately.
\end{itemize}

\vspace{\stretch{1}}

\newpage

\subsubsection*{Setting the drawing line color}

\noindent
\begin{tabularx}{\linewidth}{clcX}
\ikey{b}{blue}{Set the color of the font to blue.}
\ikey{c}{cyan}{Set the color of the font to cyan.}
\ikey{g}{green}{Set the color of the font to green.}
\ikey{k}{black}{Set the color of the font to black.}
\ikey{m}{magenta}{Set the color of the font to magenta.}
\ikey{r}{red}{Set the color of the font to red.}
\ikey{w}{white}{Set the color of the font to white.}
\ikey{y}{yellow}{Set the color of the font to yellow.}
\ikey{B}{more blue}{Increment the blue component of the current color.}
\ikey{G}{more green}{Increment the green component of the current color.}
\ikey{R}{more red}{Increment the red component of the current color.}
\ikey{$+$}{positive increment}{Set the color increment to positive.}
\ikey{$-$}{negative increment}{Set the color increment to negative.}
\end{tabularx}

\subsubsection*{Setting the drawing line size}

\noindent
\begin{tabularx}{\linewidth}{clcX}
\ikey{$>$}{increment}{Increment by one the size of the line.}
\ikey{$<$}{decrement}{Decrement by one the size of the line.}
\end{tabularx}

\vspace{\stretch{1}}

\newpage

\subsubsection*{Setting the kind of figure to draw}

In the setting mode, pressing one of the following keys enter the (still
experimental) figure drawing mode:

\noindent
\begin{tabularx}{\linewidth}{clcX}
\ikey{V}{vertical line}{Draw a vertical line.}
\ikey{H}{horizontal line}{Draw a horizontal line.}
\ikey{S}{segment}{Draw a segment.}
\ikey{C}{circle}{Draw a circle.}
\ikey{p}{point}{Draw a point.}
\ikey{P}{polygone}{Draw a polygone.}
\ikey{e}{endpoly}{Close the polygone that is beeing drawn.}
\ikey{F}{free hand}{Draw a line following the pointer.}
\ikey{' '}{cancel}{Cancel the figure setting.}
\end{tabularx}

\vspace{\stretch{1}}


\def \newpage{\Newpage}

%To make the \section*{Index} behave as \section{Index}
\let \Section \section\def \section*{\Section}
\printindex

\end{document}

% LocalWords:  DVI Didier Weis advi INRIA LGPL Miquel Furuse Botlan XBuffer pst
% LocalWords:  XSelection Hyperref hyperref dvips advirecording advianchoring
% LocalWords:  CamlImages topleft topright bottomleft bottomright hbox PStricks
% LocalWords:  advianchor eepic pstricks ps pos bg colspec bubbleset in-lined
% LocalWords:  flushpage XFlush XSync
